\section{符号と単位を読み落とさない}
グラフの高さが負になる場合，幅と高さの積も負になる。幾何学的な面積は負にはならないが，積分で足し合わせる量には符号が付く。この違いは，運動を考えると自然に理解できる。

右向きを正とする直線上で，物体が毎秒 $2$ メートルの速度で $3$ 秒進み，その後，毎秒 $1$ メートルの速度で逆向きに $2$ 秒進んだとする。位置の変化は
\[
 2\times3+(-1)\times2=4
\]
メートルである。一方，実際に動いた道のりは
\[
 2\times3+1\times2=8
\]
メートルとなる。速度のグラフで符号を付けて面積を足すと位置の変化が得られ，速さのグラフの面積を足すと道のりが得られる。同じグラフを使っていても，何を求めているかによって計算が変わる。

単位も計算の意味を教えてくれる。位置を時間で割れば速度になり，速度に時間を掛ければ位置の変化になる。面積を横の長さで割れば縦の長さになり，縦の長さに横の長さを掛ければ面積になる。微分や積分の式が複雑になっても，単位を確かめる習慣は役に立つ。

これから，まず変化率を求める方法を整え，続いて積分を定義する。その後の各章では，式をどのように変えれば既知の微分の逆を見つけられるか，あるいは不定積分を求めずに定積分の値を決められるかを調べる。最初からすべての公式を覚える必要はない。今扱っている式で，足しているのか，割っているのか，どの変数を変えているのかを確かめながら進めよう。

